\chapter{Operaciones y estructuras algebraicas}
\label{cap:estructuras}

\section{Operaciones internas}
Una \textbf{operación interna} en un conjunto $E$ es una aplicación
$E\times E\to E$, $(a,b)\mapsto a*b$: cada par tiene un único resultado
que pertenece a $E$. La suma es interna en $\N$, pero la resta no lo es.
La composición es interna en el conjunto de aplicaciones $E\to E$.

Las propiedades básicas son:
\begin{itemize}
\item \textbf{Asociatividad}: $(a*b)*c=a*(b*c)$ para todos $a,b,c\in E$.
Permite omitir paréntesis en productos de varios términos.
\item \textbf{Conmutatividad}: $a*b=b*a$ para todos $a,b\in E$.
Junto con la asociatividad permite reordenar los términos.
\item \textbf{Elemento neutro}: $e*a=a*e=a$ para todo $a\in E$.
Si existe, es único: dos neutros $e,e'$ satisfacen $e=e*e'=e'$.
\item \textbf{Simétrico de $a$}: un elemento $a'$ con $a*a'=a'*a=e$.
Si la operación es asociativa y tiene neutro, el simétrico, cuando existe,
es único y el simétrico de $a'$ es $a$.
\end{itemize}
En notación aditiva se escriben $0,-a$ y $na$; en notación multiplicativa,
$1,a^{-1}$ y $a^n$. Para operaciones asociativas, $na$ y $a^n$ expresan
repeticiones de la operación; $\sum$ y $\prod$ abrevían sumas y productos.

\section{Grupos y subgrupos}
\label{sec:grupos}
Un \textbf{grupo} $(G,*)$ tiene una operación interna asociativa, neutro
y simétrico para cada elemento. Es \textbf{abeliano} o conmutativo si,
además, $a*b=b*a$. Son abelianos $(\Z,+)$, $(\Q,+)$ y
$(\Q\setminus\{0\},\cdot)$. Las matrices invertibles de orden $n$
forman un grupo por el producto; para $n\ge2$ no es conmutativo.

En cualquier grupo:
\[
 a*b=a*c\imp b=c,\qquad b*a=c*a\imp b=c,
 \qquad (a*b)^{-1}=b^{-1}*a^{-1}.
\]
La cancelación se obtiene multiplicando por el inverso por el lado
correspondiente. Respetando el orden de los factores, las soluciones únicas son
\[
 a*x=b\eqv x=a^{-1}*b,\qquad
 x*a=b\eqv x=b*a^{-1}.
\]
Un \textbf{subgrupo} es un subconjunto $H\subseteq G$ que, con la operación
restringida, es un grupo. Un criterio que reúne todas las condiciones es
\[
 H\ne\varnothing,\qquad
 a,b\in H\imp a*b^{-1}\in H.
\]
En efecto, tomando $a=b$ se obtiene el neutro; después los inversos y,
con ellos, la estabilidad por productos. El neutro y los inversos son
los mismos que en $G$. Siempre son subgrupos $\{e\}$ y $G$.

\textbf{Clases asociadas a un subgrupo.} Si $G$ es abeliano y $H$ es un
subgrupo, la relación $a\sim b\eqv a*b^{-1}\in H$ es de equivalencia.
La clase de $a$ es $aH=\{a*h:h\in H\}$. La aplicación $h\mapsto a*h$
es una biyección de $H$ sobre $aH$, por lo que todas las clases tienen
el mismo cardinal. Si $G$ es finito y hay $k$ clases,
\[
 |G|=k|H|;
\]
en particular, $|H|$ divide a $|G|$.

\section{Anillos, subanillos e ideales}
\label{sec:anillos}
Un \textbf{anillo} $(A,+,\cdot)$ tiene dos operaciones internas:
$(A,+)$ es un grupo abeliano, el producto es asociativo y satisface
ambas leyes distributivas:
\[
 a(b+c)=ab+ac,\qquad (a+b)c=ac+bc.
\]
Es \textbf{conmutativo} si $ab=ba$. Es \textbf{unitario} si el producto
tiene un neutro $1\ne0$. La suma siempre tiene neutro $0$ y opuestos;
el producto no tiene por qué tener unidad ni inversos.
Los anillos $\Z,\Q,\R$ son conmutativos y unitarios; las matrices
cuadradas forman un anillo unitario, no conmutativo para orden $n\ge2$.

Las reglas de signos se deducen de la distributividad:
\[
 a0=0a=0,\quad (-a)b=a(-b)=-(ab),\quad (-a)(-b)=ab.
\]
Si el anillo es conmutativo,
\[
 (a+b)^2=a^2+2ab+b^2,\qquad (a+b)(a-b)=a^2-b^2.
\]
En un anillo conmutativo y unitario, el \textbf{binomio de Newton} es
\[
 (a+b)^n=\sum_{j=0}^{n}\binom nj a^{n-j}b^j\qquad(n\in\N),
 \quad a^0=b^0=1.
\]
Para un anillo conmutativo sin unidad y $n\ge2$, se escribe la misma
expansión con los extremos $a^n,b^n$ separados y la suma de $j=1$ a $n-1$,
de modo que solo intervengan potencias positivas.

Un elemento no nulo es \textbf{divisor de cero} si forma un producto
nulo con otro elemento no nulo, por alguno de los lados. Un anillo
\textbf{íntegro}, en el sentido de ausencia de divisores de cero, satisface
\[
 ab=0\imp a=0\ \text{o}\ b=0.
\]
En él se pueden cancelar factores no nulos. En un anillo unitario,
un elemento invertible no es divisor de cero. Por ejemplo, $\Z$ es
íntegro, pero un anillo de matrices de orden $n\ge2$ tiene divisores de cero.

Un \textbf{subanillo} $H\subseteq A$ se caracteriza por
\[
 H\ne\varnothing,\qquad a,b\in H\imp a-b\in H\ \text{y}\ ab\in H.
\]
No se exige que contenga la unidad de $A$. Por ejemplo, $2\Z$ es un
subanillo de $\Z$ y no tiene unidad para el producto.

En un anillo conmutativo, un \textbf{ideal} $I$ es un subconjunto no vacío
cerrado por diferencias y que absorbe productos por elementos del anillo:
\[
 a,b\in I\imp a-b\in I,\qquad a\in I, c\in A\imp ac\in I.
\]
Todo ideal es subanillo, pero el recíproco no es automático. Son ideales
$\{0\}$ y $A$. En un anillo conmutativo unitario, el ideal principal
generado por $a$ es $(a)=aA=\{ac:c\in A\}$; contiene a $a$ y está
incluido en cualquier ideal que lo contenga.

\section{Cuerpos}
\label{sec:cuerpos}
Un \textbf{cuerpo} es un anillo conmutativo unitario en el que cada
elemento no nulo tiene inverso multiplicativo. Equivalentemente,
$(K,+)$ y $(K\setminus\{0\},\cdot)$ son grupos abelianos, con $0\ne1$,
y se cumple la distributividad. $\Q$ y $\R$ son cuerpos; $\Z$ no lo es,
pues, por ejemplo, $2$ no tiene inverso entero.
En las clases módulo $p$ se opera mediante $[a]+[b]=[a+b]$ y
$[a][b]=[ab]$. Módulo $3$ se obtiene un cuerpo, con $[2]^{-1}=[2]$;
módulo $4$ no, pues $[2]\ne[0]$ y $[2][2]=[0]$.

Todo cuerpo carece de divisores de cero. Para $a\ne0$,
\[
 ab=ac\imp b=c,\qquad ax+b=0\eqv x=-a^{-1}b=-\frac ba.
\]
Un \textbf{subcuerpo} es un subconjunto que es cuerpo con las operaciones
restringidas. Basta que $H$ tenga al menos dos elementos, sea cerrado
por diferencias y que $ab^{-1}\in H$ siempre que $a,b\in H$ y $b\ne0$.

\section{Compatibilidad con el orden}
\label{sec:estructuras-ordenadas}
Un grupo abeliano aditivo está \textbf{ordenado} si su orden cumple
\[
 a\le b\imp a+c\le b+c.
\]
El cono $G_+=\{x:0\le x\}$ incluye al cero y permite recuperar el orden:
$a\le b\eqv b-a\in G_+$. Sumar desigualdades conserva el sentido y
tomar opuestos lo invierte.

Un \textbf{anillo conmutativo ordenado} tiene un orden compatible con la
suma y satisface $a,b\ge0\imp ab\ge0$. Si el orden es total,
\begin{itemize}
\item multiplicar por $c\ge0$ conserva una desigualdad; por $c\le0$ la invierte;
\item $a^2\ge0$; si hay unidad, $0<1$;
\item el valor absoluto es $|a|=a$ si $a\ge0$ y $|a|=-a$ si $a<0$.
Se cumple $|a|\ge0$, $|a|=0\eqv a=0$, $|ab|=|a||b|$ y
$|a+b|\le|a|+|b|$.
\end{itemize}
Un \textbf{cuerpo ordenado} lleva un orden total compatible con ambas
operaciones. En él, $a>0\imp a^{-1}>0$ y
\[
 0<a\le b\imp \frac1b\le\frac1a;
 \qquad a\le b<0\imp \frac1b\le\frac1a.
\]
No se puede invertir una desigualdad con extremos de signos distintos
usando estas reglas.

\section{Homomorfismos e isomorfismos}
\label{sec:homomorfismos}
Un \textbf{homomorfismo} conserva las operaciones de la estructura.
Para una operación, $f(a*b)=f(a)*'f(b)$; para anillos debe conservar
tanto la suma como el producto. Un homomorfismo biyectivo es un
\textbf{isomorfismo}: identifica estructuras con las mismas propiedades
algebraicas. Su inversa también conserva las operaciones.
La composición de homomorfismos es un homomorfismo.
Se llama \textbf{endomorfismo} si origen y destino son la misma estructura,
y \textbf{automorfismo} si, además, es un isomorfismo.

Para un homomorfismo de grupos $f:G\to G'$,
\[
 f(e)=e',\quad f(a^{-1})=f(a)^{-1},\quad
 \ker f=\{a\in G:f(a)=e'\},\quad \operatorname{Im}f=f(G).
\]
El núcleo y la imagen son subgrupos de $G$ y $G'$, respectivamente.
Más generalmente, las imágenes y preimágenes de subgrupos son subgrupos.
Además,
\[
 f\text{ inyectivo}\eqv\ker f=\{e\},\qquad
 f\text{ sobreyectivo}\eqv\operatorname{Im}f=G'.
\]
La primera equivalencia se explica porque
$f(a)=f(b)\eqv ab^{-1}\in\ker f$.

El \textbf{subgrupo generado por $a$} es
$\langle a\rangle=\{a^n:n\in\Z\}$; en notación aditiva,
$\langle a\rangle=\{na:n\in\Z\}$. Es la imagen del homomorfismo
$\Z\to G$, $n\mapsto a^n$ o $n\mapsto na$, y es el menor subgrupo que
contiene a $a$. Los exponentes negativos usan inversos y $a^0=e$.

Un homomorfismo de anillos conserva $+$ y $\cdot$; no se exige aquí
que conserve la unidad como condición adicional. Su imagen es subanillo
y, si el anillo de origen es conmutativo, su núcleo es un ideal.
Los homomorfismos entre cuerpos se entienden con esta misma convención.

Una aplicación entre conjuntos ordenados \textbf{conserva el orden} si
$a\le b\imp f(a)\le f(b)$. Para identificar órdenes mediante un
\textbf{isomorfismo de orden} se requiere una biyección con
\[
 a\le b\eqv f(a)\le f(b);
\]
es decir, ella y su inversa deben conservar el orden. En órdenes totales
basta una biyección que lo conserve.

Las inclusiones $\N\subset\Z\subset\Q\subset\R\subset\C$ expresan
identificaciones que conservan las operaciones; hasta $\R$ conservan
también los órdenes habituales.

\section{Suma y producto de cardinales}
Para cardinales $\alpha=\operatorname{card}A$ y $\beta=\operatorname{card}B$,
se definen $\alpha+\beta$ mediante la unión de representantes disjuntos
y $\alpha\beta=\operatorname{card}(A\times B)$.
El resultado no depende de los representantes elegidos.
Ambas operaciones son asociativas y conmutativas; sus neutros son $0$
y $1$, respectivamente. El producto distribuye respecto a la suma y
$0\alpha=0$. Además,
\[
 \alpha\le\beta\eqv \text{existe un cardinal }\gamma
 \text{ tal que }\alpha+\gamma=\beta.
\]
Estas operaciones no proporcionan en general una resta de cardinales
ni convierten los cardinales infinitos en un anillo.
